\documentclass[12pt]{article}
\usepackage[T1]{fontenc}
\usepackage[utf8]{inputenc}
\usepackage[brazil]{babel}
\usepackage[margin=1in]{geometry}
\usepackage{ragged2e}
\usepackage[normalem]{ulem}
\setlength{\parindent}{0pt}
\setlength{\parskip}{0.8\baselineskip}
\begin{document}
\justifying
\noindent Verifica-se que a questão objeto da controvérsia jurídica suscitada no recurso manejado pela parte recorrente esteve afetada no representativo da controvérsia - \textbf{Tema n. }\textbf{339} da TNU - (PU no processo n. PEDILEF 0000981-71.2018.4.01.3900/PA),  afetado em 16/08/2023 foi julgado (trânsito em julgado em 27/03/2025), por meio do qual foi elucidada a seguinte questão:

\noindent \textit{"Definir se, diante alterações empreendidas pelas Leis nº 12.702/2012 e 13.324/2016, que possibilitaram a incorporação da GACEN aos proventos de aposentadoria, incide contribuição previdenciária sobre a GACEN."}

\noindent Restou firmada a seguinte tese:

\noindent \textit{"As alterações empreendidas pelas Leis nº 12.702, de 07.08.2012 e nº 13.324, de 29.07.2016, não possibilitam a incidência da contribuição previdenciária sobre a Gratificação Especial de Atividade de Combate e Controle de Endemias - GACEN, por não compor a base de cálculo da contribuição previdenciária do PSS do servidor público federal, ex vi, artigo 4º da Lei nº 10.887, de 18.06.2004, associada ao inciso II, do artigo 55, da Lei nº 11.784, de 22.09.2008, sendo uma parcela remuneratória paga em decorrência de local de trabalho." (Tema 339 TNU).}

\noindent Com efeito, depreende-se da leitura do voto condutor do acórdão que o mesmo se encontra em conformidade com o decidido no  tema pela TNU.

\noindent A proposito, consta do acordao hostilizado: (...) Gratificação de Desempenho da Carreira da Previdência, da Saúde e do Trabalho - GDPST, Gratificações Por Atividades Específicas, Sistema Remuneratório e Benefícios, Servidor Público Civil, DIREITO ADMINISTRATIVO E OUTRAS MATÉRIAS DE DIREITO PÚBLICO

\noindent Nessas condições, o conteúdo do Acórdão recorrido é consentâneo com o entendimento da TNU, firmado em julgamento de recurso representativo de controvérsia, o que atrai a incidência da regra prevista pelo art. 14, III, \textit{b,} da Resolução CJF n. 586/2019:

\noindent \textit{\textbf{Art. 14.}}\textit{ Decorrido o prazo para contrarrazões, os autos serão conclusos ao magistrado responsável pelo exame preliminar de admissibilidade, que deverá, de forma sucessiva: (...) }\textit{\textbf{III }}\textit{- negar seguimento a pedido de uniformização de interpretação de lei federal interposto contra acórdão que esteja em conformidade com entendimento consolidado: }\textit{\textbf{b)}}\textit{ em }\uline{\textit{\textbf{recurso representativo de controvérsia pela Turma Nacional de Uniformização}}}\textit{ ou em pedido de uniformização de interpretação de lei dirigido ao Superior Tribunal de Justiça;(...).}

\noindent Posto isso, com arrimo no \uline{\textbf{art. 14, III, }}\uline{\textit{\textbf{b}}}\textit{,} da Resolução CJF n. 586/2019 c/c art. 1.030, I, do CPC, \textbf{NEGO SEGUIMENTO }\textbf{a}o\textbf{ PU}\textbf{.}

\noindent Intimem-se. Aguarde-se o decurso do prazo recursal. Após, certifique-se o trânsito em julgado e devolva-se ao juízo de origem.
\end{document}
